% ============================================================================
%  CAPÍTULO 2 — ANTECEDENTES
% ============================================================================
\chapter{Antecedentes}
\label{cap:marco}

\section{Sistema hospitalario público chileno y clasificaciones administrativas}

\label{sec:sistema-hospitalario}

La red hospitalaria pública chilena se organiza a través del Sistema Nacional de
Servicios de Salud (SNSS), que agrupa cerca de 192 establecimientos distribuidos
en 29 Servicios de Salud territoriales. La clasificación vigente ordena estos
hospitales en tres niveles de complejidad (alta, mediana y baja) atendiendo
a criterios estructurales: dotación de camas, disponibilidad de especialidades
médico-quirúrgicas y servicios de apoyo diagnóstico~\citep{Cid2024}. Esta lógica,
heredada de la clasificación histórica Tipo 1 a Tipo 4, determina tanto la
asignación de recursos como el rol de cada establecimiento dentro de la red
asistencial.

El rasgo central de esta clasificación es que describe infraestructura disponible,
no producción clínica efectiva. Un hospital de alta complejidad puede concentrar
su actividad en egresos obstétricos y quirúrgicos electivos, mientras otro del
mismo nivel resuelve principalmente urgencias de alta severidad con elevado índice
de casuística. Ambos comparten etiqueta administrativa, pero su operación, medida
por el perfil de egresos, la distribución diagnóstica y la complejidad de los
casos atendidos puede ser radicalmente distinta. La literatura ha documentado
esta brecha, mostrando que los criterios estructurales no capturan la
heterogeneidad funcional real de la red~\citep{Havranek2023, Cid2024}.

Esta limitación adquiere relevancia directa en el contexto de los Grupos
Relacionados por el Diagnóstico (GRD), sistema que cuantifica la complejidad de
los casos atendidos y habilita comparaciones ajustadas por la carga de enfermedad
de cada establecimiento. Para que esas comparaciones sean válidas, los
hospitales contrastados deben presentar perfiles de producción suficientemente
homogéneos. Cuando los grupos administrativos predefinidos se definen por
criterios estructurales, esta condición no está garantizada: establecimientos
clasificados en el mismo nivel pueden diferir sustancialmente en su casuística
real. Ello se evidenció en las distorsiones observadas tras la implementación
del mecanismo de pago GRD en 2020, donde la alta variabilidad interna de los
grupos predefinidos impidió una distribución equitativa de
recursos~\citep{Cid2024}. Esta investigación aborda precisamente
esa brecha: identificar una segmentación funcional basada en la producción
clínica real que permita comparaciones más homogéneas que las provistas por
la clasificación administrativa vigente.

\section{Grupos Relacionados por el Diagnóstico (GRD)}
\label{sec:grd}

\subsection{Definición y arquitectura técnica}
\label{sec:grd-definicion}

Los Grupos Relacionados por el Diagnóstico (GRD) constituyen un sistema de
clasificación que agrupa los egresos hospitalarios en categorías clínicamente
coherentes y con consumo de recursos similar. A diferencia de las clasificaciones
administrativas tradicionales, los GRD se construyen a partir de la información
clínica contenida en cada hospitalización: el diagnóstico principal, los
diagnósticos secundarios, los procedimientos realizados, la edad del paciente y
otras variables demográficas.

Los GRD tienen su origen a finales de la década de 1960 en la Universidad de
Yale en Estados Unidos, donde fueron desarrollados inicialmente como una
herramienta para evaluar la calidad asistencial y medir la actividad
hospitalaria. Su adopción se generalizó al incorporarse a los mecanismos de pago
prospectivo, y en la actualidad constituyen un estándar internacional utilizado
tanto en sistemas de salud públicos como privados~\citep{Cid2024}.

El sistema asigna a cada egreso un código GRD mediante un algoritmo denominado
grouper, que analiza la combinación de variables clínicas y
demográficas y determina el grupo al que pertenece el caso. Cada GRD posee un
peso relativo que refleja la intensidad esperada de recursos requeridos para
atender ese tipo de hospitalización, calculado a partir de promedios históricos
de duración de estancia y costos. Este peso permite cuantificar la complejidad
agregada de la producción de un hospital, concepto conocido como casuística o
\textit{case-mix}.

La agregación de los GRD de un establecimiento permite cuantificar su carga de
enfermedad y vincular la actividad clínica con los costos, habilitando
mecanismos de pago por actividad que incentivan la eficiencia. Dado que el
sistema fue diseñado para capturar la intensidad de atención del paciente
hospitalizado, su uso se concentra en establecimientos de alta y mediana
complejidad, donde se realiza una proporción importante de la actividad de
hospitalización de corta estancia. En la red pública chilena, la implementación
y disponibilidad sistemática de registros GRD se ha concentrado principalmente
en este tipo de hospitales~\citep{Cid2024}.


\subsection{Implementación en Chile}
\label{sec:grd-chile}

Chile adoptó el sistema GRD de manera progresiva a partir de la década de 2000,
con el objetivo de modernizar la gestión hospitalaria y transitar desde un modelo
de financiamiento por oferta (basado en presupuestos históricos) hacia uno
orientado a la demanda efectiva. La implementación del sistema ha sido liderada
por el Ministerio de Salud y ha requerido un esfuerzo sostenido de capacitación,
estandarización de procesos de codificación y desarrollo de infraestructura
informática~\citep{Cid2024}.

Hacia 2024, el sistema alcanzaba a 72 hospitales públicos de alta y mediana
complejidad, con una expansión proyectada a 80 establecimientos en los próximos
años. Esta cobertura representa la mayor parte de la actividad hospitalaria de
corta estancia del sector público, consolidando a los GRD como el principal
estándar de registro de la producción asistencial en la red. En el plano
presupuestario, la Ley de Presupuestos 2023 asignó al Programa 05
\textit{Financiamiento Hospitales por Grupo Relacionado de Diagnóstico} un total
cercano a 4,65 billones de pesos, distribuidos entre 68 establecimientos
identificados individualmente en la partida~\citep{DIPRES_Programa05_2023}.

Sin embargo, la transición desde el registro clínico-administrativo hacia su uso
pleno como herramienta de gestión y financiamiento no ha estado exenta de
dificultades. En 2020 se implementó un mecanismo de pago basado en GRD que
organizaba a los hospitales en grupos administrativos predefinidos, asignando
a cada grupo una tasa de pago diferenciada. La evaluación posterior evidenció
que estos grupos presentaban alta variabilidad interna: hospitales del mismo
grupo administrativo mostraban perfiles de casuística y costos muy distintos,
generando distorsiones en el financiamiento que impedían una distribución
equitativa de recursos~\citep{Cid2024}.
Como solución de contingencia, en 2023 se fijó una tasa base única para toda
la red, eliminando temporalmente la diferenciación por grupos pero dejando
pendiente el desafío de construir una segmentación más robusta y
estadísticamente válida.

Este episodio puso de manifiesto dos problemas estructurales. Primero, la
necesidad de contar con una metodología analítica rigurosa para construir una
segmentación funcional que refleje la realidad de los establecimientos.
Segundo, la importancia de controlar la calidad de los datos de entrada, dado
que la heterogeneidad en los procesos de codificación puede introducir sesgos
que distorsionen cualquier análisis comparativo basado en registros GRD.

\subsection{Limitaciones de los registros GRD}
\label{sec:grd-limitaciones}

Pese a la riqueza de la información que consolidan, los registros GRD presentan
limitaciones que deben ser consideradas al utilizarlos con fines analíticos.
La principal radica en la calidad heterogénea de la codificación clínica. \citet{Aguila2019} documentaron variabilidad en la precisión con que los
establecimientos registran diagnósticos y procedimientos según la CIE-10 y la
CIE-9-MC, lo que puede generar sesgos en la asignación de códigos GRD y, por
ende, en los indicadores derivados.

Adicionalmente, los archivos originales de egresos presentan variaciones de
estructura y codificación entre años. Un análisis de los archivos del período
2019--2024 evidenció inconsistencias en la codificación de caracteres (UTF-8 en
2019--2020, UTF-8-SIG en 2021, UTF-16 en 2022--2023 e ISO-8859-1 en 2024), así
como un cambio en el nombre de la columna de identificación del paciente: en
2024, el campo \texttt{CIP\_ENCRIPTADO} se registra bajo el nombre
\texttt{ID\_BENEFICIARIO}.

Estas limitaciones no invalidan el uso de los registros GRD, pero sí exigen la
implementación de controles estadísticos robustos antes de cualquier análisis
comparativo. La literatura advierte que la comparación entre hospitales mediante
GRD solo es estadísticamente válida si se garantizan tres condiciones: la
homogeneidad tecnológica del algoritmo agrupador, el ajuste de riesgo por
severidad y mortalidad de los pacientes, y la verificación de la calidad de la
codificación clínica~\citep{Aguila2019}. En este trabajo, las dos primeras
condiciones están satisfechas por el uso de un agrupador único y por la inclusión
de variables de severidad como atributos del vector institucional, mientras que la tercera no
puede verificarse directamente desde los registros y constituye una limitación
inherente al estudio, discutida en la Sección~\ref{sec:limitaciones-fuente}.
Por último, aun garantizando estas condiciones, factores externos como la localización
geográfica, la carga de urgencias o el perfil epidemiológico territorial pueden
influir en el desempeño sin ser capturados plenamente por las métricas
tradicionales~\citep{Havranek2023}.


\section{Clasificaciones clínicas: CIE-10 y CIE-9-MC}
\label{sec:cie}

Cada egreso hospitalario se codifica mediante dos sistemas estandarizados: la
CIE-10 para diagnósticos y la CIE-9-MC para procedimientos. Ambos sistemas
poseen una estructura jerárquica que permite agregar la información en un nivel
temático amplio. Los diagnósticos se agrupan por capítulo en la CIE-10
(22 capítulos), mientras que los procedimientos se agrupan en 18 secciones
principales de la CIE-9-MC. Esta granularidad permite distinguir perfiles
clínicos relevantes sin incurrir en la dispersión que generaría operar a nivel
de código individual.

Al calcular la proporción de egresos de un hospital en cada capítulo diagnóstico
y sección de procedimientos, se obtiene un vector que resume su casuística. Por
ejemplo, los establecimientos con alta concentración de egresos en el
Capítulo XV de la CIE-10 (embarazo, parto y puerperio) o en la Sección 13 de la
CIE-9-MC (procedimientos obstétricos) presentan un perfil predominantemente
obstétrico. De manera análoga, aquellos con predominio del Capítulo IX
(enfermedades del sistema circulatorio) o de la Sección 7 (procedimientos
cardiovasculares) son identificables como hospitales con mayor orientación
cardiológica. En contraste, los establecimientos generales presentan una
distribución más equilibrada entre múltiples capítulos diagnósticos y secciones
de procedimientos.

Esta representación vectorial se utiliza en el presente trabajo para
caracterizar funcionalmente cada establecimiento. Los capítulos de la CIE-10 y
las secciones de procedimientos de la CIE-9-MC empleadas se detallan,
respectivamente, en las Tablas~\ref{tab:cie10-capitulos} y
\ref{tab:cie9-secciones}.

\begin{table}[H]
\centering
\caption{Capítulos de la CIE-10 y su dominio clínico}
\label{tab:cie10-capitulos}
\begin{tabular}{cll}
\hline
\textbf{N°} & \textbf{Rango} & \textbf{Dominio clínico} \\
\hline
I    & A00--B99 & Ciertas enfermedades infecciosas y parasitarias \\
II   & C00--D49 & Neoplasias \\
III  & D50--D89 & Enfermedades de la sangre y órganos hematopoyéticos \\
IV   & E00--E89 & Enfermedades endocrinas, nutricionales y metabólicas \\
V    & F01--F99 & Trastornos mentales y del comportamiento \\
VI   & G00--G99 & Enfermedades del sistema nervioso \\
VII  & H00--H59 & Enfermedades del ojo y sus anexos \\
VIII & H60--H95 & Enfermedades del oído y de la apófisis mastoides \\
IX   & I00--I99 & Enfermedades del aparato circulatorio \\
X    & J00--J99 & Enfermedades del aparato respiratorio \\
XI   & K00--K95 & Enfermedades del aparato digestivo \\
XII  & L00--L99 & Enfermedades de la piel y del tejido subcutáneo \\
XIII & M00--M99 & Enfermedades del aparato musculoesquelético y tejido conectivo \\
XIV  & N00--N99 & Enfermedades del aparato genitourinario \\
XV   & O00--O9A & Embarazo, parto y puerperio \\
XVI  & P00--P96 & Ciertas afecciones originadas en el período perinatal \\
XVII & Q00--Q99 & Malformaciones congénitas, deformidades y anomalías cromosómicas \\
XVIII& R00--R99 & Síntomas, signos y resultados anormales no clasificados \\
XIX  & S00--T88 & Lesiones traumáticas, envenenamientos y causas externas \\
XX   & V00--Y99 & Causas externas de morbilidad \\
XXI  & Z00--Z99 & Factores que influyen en el estado de salud \\
XXII & U00--U99 & Códigos para situaciones especiales \\
\hline
\end{tabular}
\end{table}

\begin{table}[H]
\centering
\caption{Secciones de procedimientos de la CIE-9-MC}
\label{tab:cie9-secciones}
\begin{tabular}{cll}
\hline
\textbf{Sección} & \textbf{Rango} & \textbf{Dominio de procedimientos} \\
\hline
00  & 00--00 & Procedimientos no clasificados bajo otros conceptos \\
01  & 01--05 & Operaciones sobre el sistema nervioso \\
02  & 06--07 & Operaciones sobre el sistema endocrino \\
03  & 08--16 & Operaciones sobre el ojo \\
03A & 17--17 & Otros procedimientos diagnósticos y terapéuticos diversos \\
04  & 18--20 & Operaciones sobre el oído \\
05  & 21--29 & Operaciones sobre la nariz, boca y faringe \\
06  & 30--34 & Operaciones sobre el aparato respiratorio \\
07  & 35--39 & Operaciones sobre el aparato cardiovascular \\
08  & 40--41 & Operaciones sobre el sistema hemático y linfático \\
09  & 42--54 & Operaciones sobre el aparato digestivo \\
10  & 55--59 & Operaciones sobre el aparato urinario \\
11  & 60--64 & Operaciones sobre órganos genitales masculinos \\
12  & 65--71 & Operaciones sobre órganos genitales femeninos \\
13  & 72--75 & Procedimientos obstétricos \\
14  & 76--84 & Operaciones sobre el aparato musculoesquelético \\
15  & 85--86 & Operaciones sobre el aparato tegumentario \\
16  & 87--99 & Procedimientos diagnósticos y terapéuticos misceláneos \\
\hline
\end{tabular}
\end{table}

\section{Experiencia internacional en perfilamiento hospitalario}

\label{sec:internacional}

El uso de técnicas de agrupamiento para caracterizar hospitales según su
producción clínica se ha extendido en diversos sistemas de salud. La evidencia
disponible muestra que el análisis de datos GRD, denominados DRG en la
literatura internacional, permite identificar perfiles de producción, patrones
de movilidad de pacientes y agrupamientos de complejidad similar que las
clasificaciones administrativas tradicionales no hacen explícitos. A
continuación se revisan tres experiencias representativas. Estas corresponden
al análisis de redes en Francia, la combinación de análisis factorial y
agrupamiento en China, y el modelamiento econométrico de costos en Suiza.

\subsection{Sistema francés}

\label{sec:francia}

El estudio de \citet{Chrusciel2021} constituyó un antecedente relevante en la
aplicación de métodos de aprendizaje no supervisado para caracterizar la red
hospitalaria pública francesa. Los autores analizaron 427 hospitales públicos
que concentraban cerca de 8{,}8 millones de hospitalizaciones durante 2016,
utilizando datos del sistema nacional de pago prospectivo basado en DRG. El
agrupamiento se realizó mediante el algoritmo \textit{k}-means a partir de tres
variables. Estas fueron el número de estancias como aproximación al tamaño del
hospital, la diversidad efectiva de la actividad clínica derivada del índice de
Simpson y un indicador de movilidad que medía la proporción de estancias seguidas
por otra hospitalización en un establecimiento del mismo Grupo Hospitalario
Regional. El número óptimo de grupos se fijó en tres mediante una regla de
mayoría aplicada a 30 criterios de validación y el método del codo.

El análisis identificó tres grupos con perfiles diferenciados. El primero reunió
34 hospitales de gran tamaño, con una mediana de 82{,}108 estancias anuales y
1{,}129 camas. Estos establecimientos presentaban una actividad altamente
diversificada y operaban como centros de referencia terciaria y cuaternaria. El
segundo grupo agrupó 236 hospitales de tamaño medio, con una mediana de 258
camas. El tercer grupo comprendió 157 establecimientos pequeños, con una mediana
de 48 camas, menor diversidad de actividad y una mayor proporción de pacientes
que posteriormente acudían a otros hospitales.

El hallazgo central fue que la red pública francesa se organiza en torno a tres
categorías funcionales con patrones de movilidad de pacientes desequilibrados.
Los hospitales de tamaño medio derivaban con mayor frecuencia pacientes hacia
los grandes centros de referencia. Por su parte, los establecimientos pequeños
derivaban pacientes tanto hacia hospitales medianos como hacia grandes centros.
Los hospitales de mayor tamaño actuaban preferentemente como receptores dentro de
la red. Para este trabajo, la relevancia del estudio radica en que muestra que un
agrupamiento no supervisado, construido sobre un número reducido de variables de
producción y de red, permite identificar una estructura funcional hospitalaria
relevante para la planificación sanitaria~\citep{Chrusciel2021}.

\subsection{Sistema chino}

\label{sec:china}

La experiencia china ha sido prolífica en la aplicación de técnicas de
agrupamiento sobre datos DRG, dada la complejidad y heterogeneidad de su sistema
hospitalario. \citet{Liu2025} evaluaron el desempeño de 306 hospitales generales
de la provincia de Sichuan, correspondientes a 173 hospitales secundarios y 133
terciarios. Para ello, combinaron análisis factorial exploratorio y agrupamiento
jerárquico a partir de seis indicadores DRG. Estos fueron el índice de
casuística, el número de GRD, el peso total, el índice de eficiencia de costos,
el índice de eficiencia temporal y la mortalidad de casos de riesgo bajo y medio.
Los indicadores se organizaron según el marco de calidad de Donabedian, que
considera estructura, proceso y resultado.

El análisis factorial condensó los seis indicadores en tres factores asociados
con capacidad médica, eficiencia y seguridad. Posteriormente, el agrupamiento se
realizó por separado dentro de los hospitales secundarios y terciarios. Esta
decisión permitió evitar que las diferencias estructurales entre ambos niveles
determinaran los resultados. En cada nivel se identificaron tres grupos de
desempeño denominados ``excelente'', ``medio'' e ``inferior''. Estas
denominaciones se establecieron a partir de los valores promedio de los
indicadores DRG en cada conglomerado.

El resultado más relevante para este trabajo es que los hospitales de una misma
categoría administrativa presentaban perfiles de desempeño marcadamente
distintos. Entre los hospitales secundarios, el grupo ``inferior'' presentó el
mayor índice de casuística y, a la vez, la mayor mortalidad de casos de riesgo
bajo y medio. Además, registró menores valores en el número de GRD y en el peso
total de la actividad. Por tanto, una mayor complejidad promedio de los casos
atendidos no se tradujo necesariamente en un mejor desempeño global. El estudio
muestra que los indicadores derivados de GRD, combinados mediante análisis
factorial y agrupamiento jerárquico, permiten identificar heterogeneidad de
desempeño que la categoría administrativa por sí sola no representa
adecuadamente~\citep{Liu2025}.

A diferencia de dicho estudio, la presente investigación no busca evaluar el
desempeño, la eficiencia ni la calidad asistencial de los hospitales. Su interés
se concentra en caracterizar perfiles funcionales a partir de indicadores de
producción, complejidad y casuística derivados de los registros GRD.

\subsection{Sistema suizo}

\label{sec:suiza}

El trabajo de \citet{Havranek2023} abordó la heterogeneidad hospitalaria en
Suiza mediante un modelo estadístico de regresión orientado a explicar las
diferencias en los costos hospitalarios. A partir de datos de 116 hospitales y
más de un millón de casos correspondientes a 2019, los autores construyeron un
modelo de regresión para explicar diferencias en los costos hospitalarios
promedio ajustados por \textit{case-mix}. El modelo final incluyó cinco
atributos estructurales. Estos fueron el número de altas, la proporción de
ingresos de urgencia o ambulancia, la tasa de GRD por paciente, el potencial de
pérdida esperado según la mezcla de GRD y la ubicación en una gran aglomeración.

El modelo explicó aproximadamente la mitad de la variación observada en los
costos hospitalarios promedio ajustados por casuística. Los resultados indicaron
que una parte importante de las diferencias de costos se asociaba con factores
estructurales y regionales que los establecimientos no controlan directamente.
Además, al incorporar las cinco variables seleccionadas, las categorías de tipo
hospitalario no añadieron una mejora sustantiva en la capacidad explicativa del
modelo. La selección de variables se evaluó mediante validación cruzada entre
los años 2017 y 2019, mostrando resultados consistentes pese a las diferencias
en las muestras de hospitales disponibles en cada año.

Para este trabajo, la relevancia del estudio suizo radica en que muestra que un
conjunto reducido de indicadores estructurales y de producción puede explicar
una parte importante de la heterogeneidad entre hospitales. Asimismo, ilustra
la utilidad de procedimientos de selección de variables para construir modelos
parsimoniosos y reducir la incorporación de variables redundantes o altamente
colineales. Aunque el estudio se orienta a explicar costos y no a construir
conglomerados hospitalarios, aporta un antecedente metodológico para seleccionar
indicadores relevantes en el análisis de la heterogeneidad funcional
hospitalaria~\citep{Havranek2023}.

\subsection{Síntesis comparativa}

\label{sec:sintesis-internacional}

Los tres estudios muestran que los datos de producción hospitalaria contienen
información relevante que no queda plenamente representada por las tipologías
convencionales. El caso chino muestra directamente que hospitales de una misma
categoría administrativa pueden presentar perfiles de desempeño y producción
distintos. El estudio suizo evidencia que un conjunto reducido de variables
estructurales permite explicar parte importante de las diferencias de costos
hospitalarios. El caso francés identifica una organización funcional de red a
partir de variables de producción, diversidad clínica y movilidad de pacientes,
sin realizar una comparación explícita con una clasificación administrativa
previa.

En conjunto, estas experiencias respaldan el uso de indicadores derivados de
datos hospitalarios para caracterizar la heterogeneidad entre establecimientos.
También muestran que la selección de variables, la validación de los resultados
y la interpretación contextual de los grupos son decisiones metodológicas
centrales. La Tabla~\ref{tab:comparativa-internacional} sintetiza las
características metodológicas de los tres estudios y su aporte al presente
trabajo.

\begin{table}[H]

\centering

\caption{Síntesis comparativa de experiencias internacionales en perfilamiento hospitalario}

\label{tab:comparativa-internacional}

\footnotesize

\setlength{\tabcolsep}{3pt}

\renewcommand{\arraystretch}{1.15}

\begin{tabular}{@{}
>{\raggedright\arraybackslash}p{1.1cm}
>{\raggedright\arraybackslash}p{2.4cm}
>{\raggedright\arraybackslash}p{2.4cm}
>{\raggedright\arraybackslash}p{3.1cm}
>{\raggedright\arraybackslash}p{3.8cm}@{}}

\hline

\textbf{País} & \textbf{Datos} & \textbf{Método} &
\textbf{Variables} & \textbf{Aporte al presente trabajo} \\

\hline

Francia

  & 427 hospitales públicos y 8{,}8\,M de hospitalizaciones durante 2016

  & \textit{k}-means y validación mediante 30 criterios

  & Número de estancias, diversidad clínica y movilidad de pacientes

  & Muestra que pocas variables de producción y red permiten identificar una estructura funcional relevante para la planificación \\[5pt]

China

  & 306 hospitales de Sichuan, secundarios y terciarios

  & Análisis factorial exploratorio y agrupamiento jerárquico

  & Casuística, número de GRD, peso total, eficiencia de costos, eficiencia temporal y mortalidad

  & Muestra diferencias de desempeño y producción entre hospitales de una misma categoría administrativa \\[5pt]

Suiza

  & 116 hospitales y más de 1\,M de casos correspondientes a 2019

  & Regresión ponderada, selección de variables y validación temporal

  & Altas, ingresos de urgencia, tasa de GRD por paciente, pérdida esperada y ubicación geográfica

  & Ilustra la selección de atributos estructurales asociados a diferencias de costos hospitalarios \\

\hline

\end{tabular}

\end{table}

\section{Antecedentes nacionales en caracterización hospitalaria}
\label{sec:nacional}

\subsection{Trabajos previos del Departamento de Ingeniería Informática USACH}
\label{sec:usach}

En el contexto nacional, el problema de la caracterización funcional de
hospitales ha sido abordado por una línea de investigación desarrollada en el
Departamento de Ingeniería Informática de la Universidad de Santiago de Chile.
Esta línea de trabajo, ha explorado sistemáticamente la
aplicación de técnicas computacionales avanzadas para superar las limitaciones de
las clasificaciones administrativas vigentes.

El trabajo de \citet{Villalobos2016} introdujo un enfoque basado en teoría de grafos para construir grupos homogéneos de hospitales según su casuística, en un contexto donde los datos GRD no estaban aún plenamente consolidados. Los autores desarrollaron una metodología para identificar grupos homogéneos de establecimientos comparables, permitiendo evaluar eficiencia técnica entre hospitales con perfiles de producción equivalentes dentro de cada grupo.

\citet{Carlier2020} evaluaron directamente si la categorización
administrativa del Ministerio de Salud se corresponde con el \textit{case-mix}
real de los hospitales. Mediante un algoritmo genético que combina selección de
variables y clustering, los autores buscaron subconjuntos de indicadores del DEIS
que replicaran los cinco grupos MINSAL (alta complejidad adultos, alta complejidad
pediátrica, mediana complejidad, baja complejidad y psiquiátricos) usando datos
del período 2014--2018. Aunque se encontraron conjuntos de variables con similitud
superior a 0{,}9 respecto a la clasificación oficial, estas variables eran
procedimientos muy específicos que cambiaban entre años, sin rasgos comunes
estables a lo largo del período. El grupo de mediana complejidad fue el más
inconsistente entre años. La conclusión es que no existe evidencia suficiente para
establecer que la categorización MINSAL refleja el \textit{case-mix} de los
hospitales, por lo que los autores recomiendan construir agrupamientos alternativos
basados en variables de producción clínica para cualquier análisis de eficiencia.

\citet{Giglio2021} extendieron esta línea aplicando un enfoque
basado en CRISP-DM sobre 187 hospitales públicos chilenos con datos del período
2014--2018. A diferencia del trabajo anterior, seleccionaron variables del DEIS
filtradas por correlación con los pesos DRG, aproximando así el \textit{case-mix}
sin requerir la implementación completa del sistema. Los autores evaluaron la
categorización MINSAL mediante índices internos de validez de clustering
(Silhouette y Ball-Hall) obteniendo valores de Silhouette inferiores a 0{,}2 en
todos los años analizados, con valores negativos en algunos casos, lo que indica
que varios hospitales estaban mal asignados según el criterio de case-mix.
Compararon además particiones alternativas obtenidas mediante K-means, PAM,
CLARA, FANNY, DBSCAN y un algoritmo NSGA-II multiobjetivo, encontrando que este
último generó sistemáticamente las mejores particiones en términos de cohesión y
separación, recomendando 3 o 4 grupos como configuración óptima según los datos
de distancia euclidiana. El trabajo confirmó lo planteado por Carlier et al.:
la categorización MINSAL no es la primera recomendación para análisis de
eficiencia técnica, y propuso una nueva categorización de consenso basada en los
mejores resultados por año como punto de partida para trabajos futuros.

Los hallazgos de estas tres investigaciones muestran evidencia de que las
etiquetas administrativas vigentes no reflejan adecuadamente el
\textit{case-mix} de los hospitales y de que los métodos de agrupamiento basados
en datos de producción generan particiones con mayor coherencia interna. No
obstante, estos trabajos también evidenciaron limitaciones metodológicas que
motivan la presente investigación.

\subsection{Brechas identificadas en la evidencia nacional}
\label{sec:brechas-nacionales}

El análisis de los antecedentes nacionales permite identificar tres brechas
principales que la presente investigación busca abordar.

\textbf{Ausencia de un pipeline aplicado a datos GRD recientes con validación
estadística y detección de atípicos.} Si bien trabajos previos de la misma línea
como \citet{Giglio2021} emplearon marcos estructurados
como CRISP-DM, ninguno aplicó un flujo de extremo a extremo sobre datos GRD del
período más reciente que integre validación estadística de la solución de
agrupamiento y detección formal de establecimientos atípicos. Esta ausencia
limita tanto la reproducibilidad de los análisis como la posibilidad de
actualizar los perfiles de forma sistemática a medida que se incorporan nuevos
años de datos.

\textbf{No integración de dimensiones clínicas con indicadores de gestión.} Los
estudios nacionales previos se focalizaron en variables agregadas de producción y
complejidad (número de egresos, peso GRD medio, severidad media), pero no
incorporaron la dimensión de casuística clínica mediante las clasificaciones
CIE-10 y CIE-9-MC. Esta omisión impide distinguir perfiles funcionales
específicos como hospitales obstétricos, cardiovasculares o quirúrgicos
que son fundamentales para interpretar clínicamente los grupos y orientar la
gestión.

\textbf{Detección de atípicos limitada a un único criterio ad-hoc.} De los
trabajos revisados en esta línea, solo \citet{Villalobos2016} abordó la
detección de establecimientos atípicos, identificando hospitales divergentes
mediante el método de rango intercuartílico sobre las distancias entre
perfiles, sin evaluar la sensibilidad del resultado a la elección del
criterio ni combinarlo con otras señales de atipicidad. \citet{Carlier2020} y
\citet{Giglio2021} no incorporaron un mecanismo equivalente. La identificación de
establecimientos atípicos es crucial para garantizar comparaciones
equitativas entre establecimientos: un hospital puede no pertenecer
claramente a ningún grupo consolidado, ya sea por su perfil híbrido, por
problemas en la calidad de sus datos o por características únicas de su
territorio. Sin mecanismos que combinen múltiples criterios y evalúen su
robustez, estos casos pueden distorsionar los centroides de los grupos o
generar comparaciones injustas.

Estas tres brechas definen el espacio de contribución de la presente
investigación. Respecto a los trabajos previos, este estudio incorpora tres
elementos diferenciadores: (i) una cobertura temporal extendida con datos GRD
del período 2019--2024, lo que permite observar tendencias y estabilidad de los
perfiles (ii) la integración de la casuística clínica mediante los capítulos
CIE-10 y las secciones CIE-9-MC como dimensión central de la representación
hospitalaria, habilitando la distinción de perfiles funcionales específicos y
(iii) la combinación de múltiples criterios de detección de atípicos
(Silhouette por observación, Isolation Forest y distancia normalizada al
centroide) con evaluación explícita de sensibilidad a sus parámetros, en
lugar de un único criterio ad-hoc. Isolation Forest se aplica antes del
clustering como una prueba de influencia: se excluyen temporalmente sus marcas,
se re-estima la partición y se cuantifica el cambio resultante. Los criterios que
requieren grupos (Silhouette individual y distancia al centroide) se utilizan
posteriormente para caracterizar asignaciones ambiguas o alejadas. Esta
separación evita tratar de manera automática los perfiles funcionales poco
frecuentes como errores que deban eliminarse.

\section{Metodologías computacionales aplicadas a datos hospitalarios}
\label{sec:metodologias-comp}

\subsection{Aprendizaje no supervisado para segmentación}
\label{sec:clustering}

El problema central del perfilamiento hospitalario consiste en descubrir patrones
de similitud funcional entre establecimientos cuando no existe una clasificación
de referencia válida. Las etiquetas administrativas actuales (alta, mediana y
baja complejidad) se han mostrado insuficientes para capturar la heterogeneidad
real de la red, como se documentó en la Sección~\ref{sec:sistema-hospitalario}.
Esta situación plantea un desafío metodológico, ya que no es posible entrenar un
modelo supervisado que aprenda de etiquetas que precisamente se busca mejorar.

El aprendizaje no supervisado resuelve este problema al descubrir estructura en
los datos sin requerir etiquetas previas. A diferencia de los métodos
supervisados, que aprenden una función que mapea características a categorías
conocidas, los algoritmos de clustering infieren agrupamientos directamente de
las similitudes y diferencias observadas en los atributos medidos. Esta capacidad
de encontrar patrones sin supervisión es fundamental cuando el objetivo es
precisamente cuestionar y redefinir la taxonomía existente.

En el contexto hospitalario, el aprendizaje no supervisado permite que los propios
datos clínicos, reflejados en variables como casuística, complejidad del
case-mix, estancia media y perfiles de procedimientos, revelen qué
hospitales funcionan de manera similar, independientemente de su etiqueta
administrativa. Esta aproximación ha sido validada internacionalmente en sistemas
de salud complejos, donde ha demostrado revelar estructuras funcionales de la red
que las clasificaciones administrativas no hacen
explícitas~\citep{Chrusciel2021, Liu2025}.

Para datos tabulares con tamaño muestral moderado, las técnicas de clustering
más utilizadas en la literatura de caracterización hospitalaria son el clustering
jerárquico y el particional (K-means y variantes). Estas técnicas han demostrado
eficacia en contextos similares: \citet{Chrusciel2021} aplicaron
K-means sobre 427 hospitales en Francia, y \citet{Liu2025} emplearon
clustering jerárquico aglomerativo sobre 306 hospitales en China. El clustering
basado en densidad (DBSCAN) puede producir resultados inestables por su
sensibilidad a los hiperparámetros en muestras de este orden, y el clustering
espectral requiere definir una función de kernel y un número de vectores propios
a retener, decisiones adicionales que introducen una complejidad de
interpretación no justificada cuando no hay evidencia previa de que la red
hospitalaria presente una geometría no convexa que este método esté
especialmente capacitado para capturar.

El clustering jerárquico aglomerativo construye una jerarquía de grupos uniendo
iterativamente los pares de observaciones o clusters más cercanos según una
medida de distancia. Esta aproximación presenta dos ventajas específicas para el
problema hospitalario. Primero, no requiere especificar a priori el
número de grupos, permitiendo explorar distintos niveles de granularidad mediante
un dendrograma, propiedad valiosa cuando no existe consenso institucional sobre
cuántos perfiles funcionales deberían existir. Segundo, el método de Ward, que
minimiza el incremento de varianza intra-grupo en cada fusión, produce grupos
compactos y equilibrados, propiedades deseables cuando se busca maximizar la
homogeneidad funcional dentro de cada perfil.

El algoritmo K-means y su variante K-means++ constituyen enfoques particionales
complementarios. K-means asigna cada observación al centroide más cercano
mediante un proceso iterativo de reasignación y recalculo, minimizando la suma de
distancias cuadradas intra-grupo. K-means++ mejora la inicialización de
centroides mediante un muestreo probabilístico que favorece puntos distantes,
reduciendo la sensibilidad a condiciones iniciales aleatorias. Su eficacia en
caracterización hospitalaria está respaldada por el trabajo de \citet{Chrusciel2021}, donde K-means sobre tres variables de producción y red
produjo una segmentación interpretable y estable de 427 hospitales públicos
franceses.

La ausencia de una clasificación funcional validada plantea el desafío de evaluar
si una solución de clustering es buena cuando no existe una verdad absoluta con
la cual compararla. La literatura ha desarrollado métricas internas de calidad
que evalúan propiedades deseables de los grupos sin requerir etiquetas de
referencia.

El coeficiente de Silhouette~\citep{Rousseeuw1987} cuantifica, para cada
observación, qué tan bien asignada está a su grupo comparando la distancia media
intra-grupo con la distancia media al grupo vecino más cercano. El coeficiente
toma valores en el rango $[-1, 1]$ donde valores cercanos a 1 indican que la
observación está más cerca de los miembros de su propio grupo que de
los de otros grupos, mientras que valores negativos sugieren asignaciones
ambiguas. El promedio del coeficiente sobre todas las observaciones proporciona
una medida global de la calidad de la partición.

El índice de Calinski-Harabasz cuantifica la relación entre la dispersión
inter-grupos y la dispersión intra-grupos. Una partición óptima maximiza la
separación entre grupos mientras minimiza la heterogeneidad dentro de cada grupo.
Valores altos de este índice favorecen soluciones con grupos compactos y bien
separados.

El índice de Davies-Bouldin mide la peor relación de similitud entre cada grupo
y su vecino más cercano a diferencia del anterior, valores bajos indican mejor
separación. Esta métrica penaliza soluciones donde existen pares de grupos con
alta similitud interna pero baja separación mutua, situación indeseable desde la
perspectiva de interpretabilidad~\citep{Arbelaitz2013}.

Más allá de las métricas puntuales, la estabilidad de las soluciones constituye
un criterio de validación complementario. Técnicas de remuestreo como el
bootstrapping evalúan la robustez de las asignaciones frente a
variaciones en los datos de entrada: se generan múltiples réplicas del conjunto
de datos mediante muestreo con reemplazo y se aplica el algoritmo en cada
réplica, midiendo luego la consistencia de las asignaciones. Soluciones que se
fragmentan ante pequeñas perturbaciones revelan una estructura frágil sin valor
interpretativo.

\subsection{Reducción de dimensionalidad y selección de variables}
\label{sec:reduccion}

Los registros hospitalarios se caracterizan por su alta dimensionalidad: cada
egreso contiene decenas de variables clínicas, demográficas y administrativas.
Al agregar esta información a nivel de hospital, construyendo por ejemplo
vectores de casuística mediante los capítulos CIE-10 y las secciones
principales CIE-9-MC más indicadores de gestión como estancia media, índice de
severidad y peso GRD medio, la dimensionalidad puede alcanzar o superar las 60
variables.

Cuando el número de características es del mismo orden que el número de
observaciones, se presentan tres problemas fundamentales. Primero, la
\textit{maldición de la dimensionalidad}, ya que en espacios de alta dimensión
las distancias euclidianas entre puntos tienden a concentrarse, lo que dificulta
la discriminación entre observaciones cercanas y lejanas y degrada el desempeño
de algoritmos de clustering basados en distancias. Segundo, el riesgo de capturar
ruido en lugar de patrones generalizables, lo que se manifiesta en soluciones
inestables donde pequeñas perturbaciones producen asignaciones radicalmente
distintas. Tercero, la pérdida de interpretabilidad, pues caracterizar un perfil
hospitalario mediante decenas de variables dificulta la comunicación de los
resultados y limita la utilidad práctica para la toma de decisiones.

Existen dos paradigmas para enfrentar la alta dimensionalidad: la reducción
mediante transformación y la selección de variables.

El análisis de componentes principales (PCA) transforma el conjunto original de
variables en un nuevo conjunto de componentes ortogonales, combinaciones lineales
de las variables originales, ordenados por la cantidad de varianza que explican.
Al retener únicamente los primeros componentes se condensa la información en un
espacio de menor dimensión preservando la mayor parte de la variabilidad de los
datos. Su limitación principal es que los componentes resultantes son
combinaciones de todas las variables originales, lo que dificulta la
interpretación clínica directa de cada dimensión.

Las técnicas de selección de variables, por el contrario, identifican el
subconjunto de atributos más relevantes sin transformarlos, preservando su
significado original. Métodos como LASSO (\textit{Least Absolute Shrinkage and
Selection Operator})~\citep{Tibshirani1996} y Elastic Net~\citep{Zou2005} aplican
penalizaciones que fuerzan a cero los coeficientes de las variables irrelevantes,
efectuando selección automática. LASSO induce escasez, de modo que solo un subconjunto
pequeño de variables recibe coeficientes no nulos, facilitando la interpretación.
Elastic Net combina las penalizaciones $L_1$ (LASSO) y $L_2$ (Ridge), permitiendo
retener grupos de variables correlacionadas cuando esto es apropiado.

\citet{Bollmann2023} demostraron, en el contexto de datos de
pacientes anidados en hospitales suizos, que incorporar la estructura
jerárquica del dato al aplicar LASSO produce una selección de variables
menos sesgada que ignorar las diferencias estructurales entre
establecimientos: variables que aparentaban ser altamente relevantes al
ignorar el hospital perdían por completo su importancia al considerar los
efectos de cada establecimiento, revelando que su aparente relevancia
reflejaba diferencias entre hospitales y no características de los
pacientes. En el contexto de este trabajo, la selección de variables
permite identificar qué capítulos CIE-10, qué secciones CIE-9-MC y qué
indicadores de gestión son verdaderamente discriminantes para caracterizar los
perfiles hospitalarios, descartando atributos que aportan ruido sin valor
informativo.

\subsection{Detección de anomalías}
\label{sec:anomalias}

La identificación de hospitales con comportamiento divergente es un componente
esencial del perfilamiento funcional, no una complicación técnica a ignorar. Un
establecimiento puede presentar un patrón atípico por razones legítimas que no
deben forzarse dentro de un grupo convencional.

Una primera causa son los perfiles híbridos o de transición. Un hospital que
actúa como centro de referencia regional para especialidades complejas y a la
vez como prestador principal de atención general para su zona de influencia
directa puede exhibir una mezcla de casuística que no se ajusta a los perfiles
típicos de alta ni de mediana complejidad. \citet{Miyata2010}
documentaron este fenómeno al aplicar un modelo de predicción de mortalidad
ajustado por casuística sobre datos de 469 hospitales japoneses: el modelo
discriminaba bien en la generalidad de los establecimientos, pero su ajuste
era significativamente peor en el subconjunto de hospitales especializados
y en aquellos con salas de convalecencia de larga estancia, precisamente los
que combinan un perfil funcional mixto que no encaja en el patrón de
hospital agudo general usado para construir el modelo.

Una segunda causa son los problemas de codificación clínica. Como se documentó
en la Sección~\ref{sec:grd-limitaciones}, la heterogeneidad en los procesos de
registro puede generar datos anómalos que no reflejan diferencias funcionales
reales sino artefactos del sistema~\citep{Aguila2019}. Un hospital con baja
calidad de codificación puede aparecer como atípico simplemente porque sus datos
no son comparables con los del resto de la red.

Una tercera causa son las características singulares del contexto territorial o
histórico. Establecimientos ubicados en zonas con perfiles epidemiológicos
excepcionales, o con funciones históricas que transitaron desde la
especialización hacia la atención general, pueden presentar patrones que no se
replican en ningún otro establecimiento. La evidencia internacional respalda esta
idea, pues \citet{Havranek2023} mostraron que factores regionales como
la ubicación geográfica explican variaciones sistemáticas entre hospitales que
persisten aun después de ajustar por casuística.

Identificar y caracterizar estos casos es fundamental por dos motivos. Incluir un
hospital atípico en un grupo convencional distorsiona el centroide del grupo y
genera comparaciones injustas tanto para el establecimiento divergente como para
los miembros típicos. Además, los hospitales atípicos no son errores del sistema
sino parte de la complejidad que una clasificación funcional debe reconocer de
manera explícita.

El algoritmo LOF (\textit{Local Outlier Factor})~\citep{Breunig2000} detecta
atípicos basándose en la densidad local de las observaciones. A diferencia de
métodos globales que definen un outlier como un punto lejano al centroide
general, LOF considera la estructura local, de modo que un punto es considerado
anómalo si su vecindad es significativamente menos densa que la de sus vecinos. Esta
aproximación resulta robusta para identificar atípicos en contextos donde existen
múltiples grupos de densidades heterogéneas, como ocurre en datos hospitalarios
cuando coexisten perfiles de distintas escalas y complejidades.

La ventaja de LOF radica en su capacidad de detectar anomalías contextuales, ya
que un hospital puede no ser atípico en el espacio global, pero sí dentro de su
grupo asignado. Por ejemplo, un establecimiento clasificado en un clúster de hospitales
generales de mediana complejidad pero con un perfil inusualmente especializado en
cirugía cardiovascular puede no ser un outlier global, porque existen hospitales
cardiovasculares en otros grupos, pero sí puede ser un outlier local dentro de su grupo.

La distancia de Mahalanobis constituye una alternativa complementaria, midiendo la
distancia de cada observación al centroide del conjunto considerando la estructura
de covarianza de los datos. A diferencia de la distancia euclidiana, que trata
todas las direcciones por igual, la distancia de Mahalanobis es invariante ante
transformaciones lineales y penaliza desviaciones en direcciones de baja
varianza. Esta propiedad la hace apropiada para detectar atípicos multivariados
en datos correlacionados, como los vectores de casuística hospitalaria, donde
muchas variables están naturalmente correlacionadas.

Ambas técnicas capturan formas distintas de atipicidad. LOF identifica hospitales
con vecindades locales anómalas dentro de su grupo asignado, siendo sensible a
la estructura local de los datos. Mahalanobis detecta establecimientos alejados
del centroide considerando la correlación entre variables, siendo más adecuado
cuando las variables presentan alta covarianza. Sin embargo, ambas presentan
limitaciones operativas que deben considerarse según el contexto, ya que LOF es
sensible a la elección del parámetro de vecindad y puede producir resultados
inestables en muestras pequeñas, mientras que Mahalanobis requiere estimar la
matriz de covarianza, operación mal condicionada cuando el número de variables
es del mismo orden que el número de observaciones.

Una alternativa que sortea ambas limitaciones es el algoritmo Isolation
Forest~\citep{Liu2008}, que detecta atípicos basándose en la facilidad de aislar
una observación mediante particiones aleatorias del espacio de características.
Los puntos que se aíslan con pocas particiones son considerados atípicos. A
diferencia de LOF, no requiere definir un radio de vecindad y a diferencia de
Mahalanobis, no requiere estimar la matriz de covarianza. Estas propiedades lo
hacen especialmente adecuado para conjuntos de datos con tamaño muestral
moderado y alta dimensionalidad relativa.

La elección entre estas técnicas depende de las características del problema:
el tamaño muestral, la dimensionalidad de los datos y el tipo de atipicidad
que se busca detectar. En cualquier caso, la combinación de múltiples criterios
complementarios siendo cada uno sensible a una faceta distinta de la anomalía
produce una caracterización más robusta que la aplicación de un único detector.

\subsection{Validación y evaluación de resultados}
\label{sec:validacion}

La validación de las soluciones de clustering en ausencia de una verdad absoluta
constituye uno de los desafíos centrales del perfilamiento hospitalario. A
diferencia del aprendizaje supervisado, donde se dispone de etiquetas correctas
para medir el desempeño, en clustering no existe una partición de referencia que
defina qué es una buena solución. El problema es especialmente relevante, ya que el objetivo es precisamente superar unas clasificaciones
administrativas que han demostrado no ser representativas de la realidad
funcional. La literatura aborda este desafío mediante cuatro estrategias
complementarias:

\begin{itemize}

  \item \textbf{Validación interna.} Las métricas internas evalúan propiedades
        geométricas deseables de los grupos sin requerir etiquetas de referencia.
        Una partición de calidad debe combinar alta cohesión interna con alta
        separación entre grupos distintos, aunque ambos objetivos pueden entrar en
        tensión según la estructura de los datos. Los tres índices descritos en la
        Sección~\ref{sec:clustering} cumplen funciones complementarias: el
        coeficiente de Silhouette permite diagnosticar grupos mal cohesionados y
        orientar decisiones de refinamiento, el índice de Calinski-Harabasz es
        útil para comparar soluciones con distinto número de grupos y apoyar la
        elección de $k$ y el índice de Davies-Bouldin ayuda a detectar
        sobresegmentación, penalizando pares de grupos con centroides próximos y
        alta dispersión interna. Emplearlos de forma conjunta reduce el riesgo de
        sobreajustar a un único criterio.

  \item \textbf{Validación de estabilidad.} Una solución estable mantiene sus
        grupos principales cuando se introducen perturbaciones moderadas en los
        datos de entrada. El bootstrapping evalúa esta propiedad
        generando múltiples réplicas mediante muestreo con reemplazo, aplicando el
        algoritmo en cada réplica y midiendo la consistencia de las asignaciones
        con índices como el de Jaccard o la medida F. Un grupo inestable que se
        fragmenta o fusiona ante pequeñas variaciones revela un perfil frágil que
        requiere análisis adicional antes de fundamentar decisiones de gestión.
        La validación por estabilidad cobra mayor relevancia cuando el tamaño
        muestral es moderado, ya que en esas condiciones una sola observación
        influyente puede alterar la configuración de los grupos.

  \item \textbf{Validación externa.} Cuando existe una clasificación externa,
        como la administrativa vigente, es posible contrastar la solución empírica
        con esa referencia mediante índices de concordancia como el de Jaccard, la
        medida F o el índice de Rand ajustado~\citep{Halkidi2001}. Un bajo acuerdo
        entre la clasificación empírica y la administrativa no invalida la primera,
        de hecho, dado que la clasificación administrativa ha mostrado limitaciones
        documentadas en las Secciones~\ref{sec:sistema-hospitalario}
        y~\ref{sec:grd-chile}, es esperable que una segmentación basada en datos
        clínicos reales diverja de las etiquetas estructurales. El contraste
        permite cuantificar esa divergencia y examinar los casos donde la
        discrepancia es mayor.

  \item \textbf{Validación cualitativa por interpretabilidad clínica.}
        \citet{Wick2022} combinaron análisis estadístico con juicio
        clínico experto para segmentar usuarios de alto costo, argumentando que el
        análisis aislado puede producir grupos sin sentido clínico y que la
        revisión por expertos asegura que los subgrupos resultantes sean relevantes
        y accionables. En el contexto hospitalario, esto implica verificar que los
        perfiles emergentes sean consistentes con el conocimiento sobre
        especialización funcional y configuración de la red. Un grupo con alta
        proporción de egresos del Capítulo XV de la CIE-10 debería estar
        constituido por hospitales con unidades obstétricas robustas y uno con altos
        índices de severidad y peso GRD medio debería contener hospitales con
        capacidad de cuidados intensivos. Esta validación no reemplaza a las
        métricas cuantitativas, pero aporta un filtro de plausibilidad clínica
        indispensable para que los resultados sean adoptados por los tomadores de
        decisiones.

\end{itemize}

\subsection{Caracterización e interpretabilidad de los perfiles}
\label{sec:interpretabilidad}

Una vez asignadas las etiquetas de clúster, la caracterización de los perfiles se
apoya en técnicas estadísticas directamente interpretables. Las pruebas de
contraste por variable: Kruskal-Wallis para atributos numéricos y chi-cuadrado
para categóricos permiten identificar qué indicadores difieren de forma
significativa entre grupos y, por tanto, contribuyen a definir cada perfil. La
regresión logística regularizada con penalización LASSO modela la pertenencia a
cada grupo en función de las variables originales: sus coeficientes ofrecen una
lectura directa de los determinantes de cada perfil y, al forzar a cero los
coeficientes de las variables irrelevantes, producen una caracterización compacta
y comunicable.

Este trabajo privilegia métodos intrínsecamente interpretables en línea con el
argumento de \citet{Rudin2019}, quien sostiene que en contextos de toma de
decisiones de alta responsabilidad un modelo directamente legible es preferible
a la combinación de una caja negra con una explicación post-hoc que solo
aproxima el razonamiento subyacente. Las técnicas de explicabilidad post-hoc como los valores
SHAP~\citep{Lundberg2017} no se aplican en este trabajo dado que no hay un modelo
de base complejo que requiera explicación auxiliar, se mencionan únicamente como
posible extensión futura.

\section{Marco metodológico: CRISP-DM}
\label{sec:crisp-dm}

El desarrollo de este trabajo se fundamenta en el marco metodológico CRISP-DM
(\textit{Cross-Industry Standard Process for Data Mining})~\citep{Chapman2000}, un estándar
consolidado para proyectos de minería de datos que estructura el proceso en seis
fases iterativas: comprensión del negocio, comprensión de los datos, preparación
de los datos, modelado, evaluación y despliegue.

La elección de CRISP-DM frente a alternativas como SEMMA, KDD o CPMAI responde a
características específicas del problema abordado. SEMMA, desarrollado por SAS,
orienta su ciclo de vida casi exclusivamente hacia el procesamiento técnico de la
información bajo ecosistemas de software cerrados, con menor énfasis en la
comprensión del dominio. KDD, por su parte, mantiene un enfoque informático
centrado en la extracción y transformación, omitiendo una fase metodológica
formal para entender el problema desde su origen empírico. CPMAI, finalmente, está
estructurado para el desarrollo, despliegue y mantenimiento de soluciones
cognitivas en entornos operativos, lo que resulta desproporcionado para un
proyecto de caracterización funcional con fines de investigación.

CRISP-DM, en cambio, otorga un peso fundamental a la fase inicial de comprensión
del entorno, lo que permite traducir la complejidad clínica de los registros
hospitalarios y del sistema GRD en un problema de minería de datos claro y
abordable. Su naturaleza altamente iterativa se alinea con el objetivo de
caracterización funcional y selección de características, y resulta
especialmente apropiado cuando no existe un etiquetado previo que guíe el
modelado, pues obliga a iterar entre la comprensión de los datos y el análisis
hasta asegurar que los patrones descubiertos tengan relevancia para la gestión
sanitaria. La aplicación concreta de cada fase al problema hospitalario se
describe en el Capítulo~\ref{cap:metodologia}.
